\section{The Past}\label{sec:the-past}
In 1915, Einstein presented the theory of general relativity.
Incorporating his special relativistic view of the spacetime with the principle of equivalence (his happiest thought), he gave a geometrical description of spacetime, and gravitation as a curvature (in opposed to viewing it as a force between massive bodies), in which a freely falling body is not being pushed, but moves in the path of geodesics.
Perhaps

\begin{quote}
    ``Spacetime tells matter how to move; matter tells spacetime how to curve.''\cite{misner_gravitation_2017}
\end{quote}

explains the essence of general relativity.
The theory, despite its elegant mathematics and clean philosophy and principles, went on a slow path of earning its pride of place in physics.
Though it explained the anomalous precession of Mercury's perihelion; a first golden opportunity arose on 29 May 1919, when a total solar eclipse allowed the measurement of the deflection of light by the sun\cite{dyson_ix_1920}.

The observation were of poor quality, despite which, made Einstein, and his theory, a celebrity overnight.
Later, on 22 September 1922, another set of observations, with outstanding quality, came out with precise agreement with General Relativity's predictions\cite{campbell_observations_1923}.
From then on, additional tests has followed, first with the observation of gravitational redshift in astronomical measurements in 1954\cite{popper_red_1954}, and in laboratory in 1959\cite{pound_apparent_1960}, and later with the observation of gravitational time dilation in 1964\cite{shapiro_fourth_1964}.

However, the most spectacular prediction of general relativity, gravitational waves, remained elusive for decades.

In his 1916 paper, Einstein predicted the existence of gravitational waves\cite{einstein_grundlage_1922}, emerging from Einstein's general theory of relativity.
Just months after he formulated the field equations, in the linearized weak-field limit, he found wave-like solutions that propagate at the speed of light.
However, the physical reality of these waves remained controversial for decades.
Even Einstein himself wavered, publishing a paper in 1946 claiming gravitational waves didn't exist.

The debate was finally settled in the 1950s through the work of Bondi, Pirani, and Robinson, who demonstrated that gravitational waves could carry energy and produce observable effects on test particles\cite{bondi_gravitational_1959,bondi_plane_1957,hill_how_2016}.

The theoretical understanding of gravitational radiation accelerated in the 1960s.
The quadrupole formula, which describes the leading-order from slowly moving sources, provided the foundation for calculating energy loss from orbiting systems.
Post-Newtonian approximations extended these calculations to higher orders, allowing predictions for how binary systems should evolve as they radiate energy.
By the early 1970s, the stage was set for an indirect confirmation of gravitational waves through astrophysical observations.

\subsection{The Dance of the Pulsars}\label{subsec:the-dance-of-the-pulsars}

The 1974 discovery of the binary pulsar PSR B1913+16 by Hulse and Taylor provided the first indirect evidence for gravitational radiation.
This system—two neutron stars in a tight orbit—offered a natural laboratory for testing general relativity's predictions\cite{hulse_discovery_1975}.
As the system radiates gravitational waves, it should lose orbital energy, causing the orbital period to decrease at a precisely calculable rate.
Over the following years, radio timing observations tracked this orbital decay with extraordinary precision\cite{miller_new_2019}.
By 1993, when Hulse and Taylor received the Nobel Prize, the observed decay matched general relativity's predictions to better than 0.5\%.
Gravitational waves were real; they just hadn't been directly detected\cite{damour_1974_2015}.

\subsection{Building the Ears}\label{subsec:building-the-ears}
The quest for direct detection of gravitational waves began in earnest in the 1960s with Joseph Weber's pioneering work on resonant bar detectors.
Weber's bars were massive aluminum cylinders designed to resonate in response to passing gravitational waves.
Though Weber claimed detections in the late 1960s, his results were not reproducible, and the scientific community remained skeptical\cite{weber_evidence_1969,weber_gravitational_1967}.

The experimental work of Weber in 1969, together with his unverified claims of detection, spurred the development of more sensitive detectors.
In the 1970s and 1980s, several groups around the world began constructing interferometric detectors, which promised orders of magnitude greater sensitivity than resonant bars.
The Laser Interferometer Gravitational-Wave Observatory (LIGO) project, initiated in the late 1980s, aimed to build large-scale interferometers capable of detecting gravitational waves from astrophysical sources\cite{weiss_republication_2022}.

The idea lies behind the fact that given a passing gravitational radiation, spacetime is curved and therefore, lengths change.
Using a pair of perpendicular (for getting the maximum length change) tunnels and two beams of in-phase photon lasers (so that the speed is constant); one can observe interference fringes produced when the light is combined at the corner station shift back and forth as the wave changes in phase.
This shift can be compared with the expectation from different types of signals to assess the probability that signal or noise is being observed\cite{miller_new_2019}.

The LIGO (Laser Interferometer Gravitational-Wave Observatory) project was proposed in the 1980s and approved by the U.S.\ Congress in 1992.
Construction of two 4-kilometer facilities—one in Hanford, Washington, and one in Livingston, Louisiana—was completed in 1999.
Initial LIGO operated from 2002 to 2010 without detecting gravitational waves, as expected; the technology wasn't yet sensitive enough.
The crucial upgrade to Advanced LIGO, incorporating decades of technical improvements, was completed in 2015.


Europe built the Virgo detector near Pisa, Italy, with 3-kilometer arms, which began operations in 2007 and underwent similar sensitivity upgrades.
Japan constructed KAGRA, an underground 3-kilometer detector in the Kamioka mine, which joined the network in 2020.
This global network of detectors is essential for localizing sources on the sky and testing the polarization properties of gravitational waves.

\subsection{The New Astronomy}\label{subsec:the-new-astronomy}

Four decades of thought, development, went on so in by the end of summer of 2015 the LIGO detectors were finally sensitive enough to detect plausible astrophysical sources.
And on 14 September 2015 timing favored physicists again.
Almost as soon as the detectors were turned on, they gave a signal that was strong enough to be an unmistakable source\cite{abbott_observation_2016a}.

The signal—designated GW150914—swept through LIGO's sensitive frequency band in a tenth of a second, matching theoretical predictions from numerical relativity simulations with remarkable fidelity.
The waveform's shape encoded the masses and spins of the black holes, the distance to the system, and even confirmed that the merged remnant obeyed general relativity's predictions for black hole ringdown.
This single detection accomplished multiple scientific goals simultaneously: it confirmed gravitational waves exist, proved that binary black holes exist and merge within a Hubble time, and tested general relativity in the strong-field regime for the first time.

The announcement on 11 February 2016 marked the dawn of gravitational wave astronomy.
The 2017 Nobel Prize in Physics was awarded to Rainer Weiss, Kip Thorne, and Barry Barish for their decisive contributions to the LIGO detector and the observation of gravitational waves.

But the story was just beginning.
The truly transformative event came on August 17, 2017, with the detection of GW170817—gravitational waves from a binary neutron star merger\cite{collaboration_gw170817_2017}.
Within seconds of the gravitational wave detection, the Fermi and INTEGRAL gamma-ray satellites detected a short gamma-ray burst from the same sky region.
Within hours, optical telescopes located the host galaxy (NGC 4993) and observed the kilonova—the radioactive glow from newly synthesized heavy elements\cite{goldstein_ordinary_2017}.
Over the following weeks, telescopes across the electromagnetic spectrum, from radio to X-rays, tracked the event's evolution.
This single event spawned over 100 papers in the first few months, addressing questions ranging from the neutron star equation of state to the cosmic origin of gold and platinum, from the Hubble constant to tests of fundamental physics.

GW170817 inaugurated the era of multi-messenger astronomy with gravitational waves.
It demonstrated that gravitational wave observations could guide electromagnetic follow-up, that joint observations could constrain astrophysics and cosmology in ways neither could alone, and that the gravitational wave community could coordinate with observers worldwide to maximize scientific return.
The event also validated decades of theoretical work on neutron star mergers and showed that these systems are indeed factories for heavy element production.

\section{The Present}\label{sec:the-present}
The era of ground-based gravitational-wave detection with LIGO and Virgo is ongoing, with further instrumental improvements already planned.
One such upgrade is known as $A+$~\cite{cooper_sensors_2023}, which aims to improve broad band sensitivity through the use of quantum light squeezing.
This technique reduces phase noise at high frequencies and radiation-pressure noise at low frequencies, thereby enhancing the detectors' ability to observe a wider range of sources\cite{miller_new_2019}.

Gravitational-wave astronomy is not limited to the LIGO and Virgo observatories.
Japan has constructed a $3$-km underground interferometer, KAGRA, located in the Kamioka mine\cite{_kagra_,ooba_kagra_,akutsu_kagra_2019}.
This detector represents the first implementation of cryogenic interferometric detection.
KAGRA is intended to operate as part of the global gravitational-wave detector network alongside LIGO and Virgo\cite{ohashimasatake_ligo_}.
In addition, the LIGO-India detector is expected to become operational in the mid-2020s.
With detectors distributed globally, the network will achieve improved sensitivity to gravitational waves arriving from different directions and will operate closer to optimal performance\cite{miller_new_2019}.

Following GW170817~\cite{abbott_properties_2019}, the pace of discoveries accelerated.
The first two observing runs (O1 and O2) detected 11 gravitational wave events~\cite{abbott_gwtc1_2019}.
The third observing run (O3), split into O3a and O3b from April 2019 to March 2020, detected approximately 80 additional events.
The third Gravitational-Wave Transient Catalog (GWTC-3), released in 2021, included 90 confident detections~\cite{abbott_gwtc3_2023}.
These observations have revealed:
\begin{itemize}
    \item \textbf{A diverse population of binary black holes}, ranging from stellar-mass systems to intermediate-mass mergers, with a wide range of mass ratios and spin configurations.
    The population properties provide constraints on formation channels, stellar evolution, and binary dynamics.


    \item \textbf{Binary neutron star systems** beyond GW170817}, though none with confirmed electromagnetic counterparts.
    These detections constrain the neutron star mass distribution and equation of state.

    \item \textbf{Neutron star-black hole binaries}, confirming that this channel contributes to the merger rate and enabling new tests of compact object physics.

    \item \textbf{Evidence for hierarchical mergers}--black holes that themselves resulted from previous mergers—suggesting dense stellar environments where repeated interactions can build up increasingly massive black holes.

    \item \textbf{Tests of general relativity} with increasing precision.
    So far, all observations are consistent with GR predictions, placing stringent constraints on alternative theories of gravity.

\end{itemize}


The field has matured to the point where multiple detections per week are routine during observing runs.
The fourth observing run ($O4$) began in May 2023 with improved sensitivity across the network.
The Japanese KAGRA detector joined at reduced sensitivity, working toward full integration.
Plans for $O5$, potentially beginning in 2026--2027, promise further sensitivity improvements and higher detection rates.

\section{Looking Ahead}\label{sec:looking-ahead}
Although predictions beyond 2025 are uncertain, plans for more advanced ground-based detectors are already under discussion.
Two major proposals are the Einstein Telescope and the Cosmic Explorer.
Both concepts require new facilities.
The Einstein Telescope is planned as an underground triangular interferometer to mitigate seismic noise~\cite{etscienceteam_einstein_2011}, while the Cosmic Explorer envisions a detector with arm lengths of approximately  $40$km, incorporating quantum squeezing techniques similar to those of $A+$ and cryogenic cooling technologies developed for KAGRA and Voyager~\cite{hall_cosmic_2022}.
At present, funding for these projects has not been secured, and only early-phase studies have begun.
Consequently, deployment is not expected before the 2030s.


In addition to the frequency range of approximately $10\textendash2,000 \text{Hz}$ probed by current and planned ground-based detectors, lower-frequency gravitational waves can be studied indirectly through cosmic phenomena.
These approaches provide access to gravitational-wave signals outside the reach of terrestrial interferometers.

\subsection{Space-Based Detectors}\label{subsec:space-based-detectors}


To detect gravitational waves emitted by different classes of strong-gravity sources, space-based detectors are required.
One such mission is the Laser Interferometer Space Antenna (LISA), a mission led by the European Space Agency with participation from NASA. LISA will consist of three spacecraft arranged in a triangular formation, orbiting the Sun while trailing the Earth and maintaining a precessing configuration.


LISA is expected to detect new classes of gravitational-wave sources, including mergers of black holes with masses in the range $10^4 M_\odot$ to $10^7 M_\odot$ located at galactic centers.
The gravitational waves emitted by these systems contain information about the properties of black holes prior to and during galaxy mergers, contributing to astrophysical population studies.


Space-based gravitational-wave observations are particularly important due to the unique science they enable.
Information about supermassive black holes during galaxy mergers can only be obtained through these observations.
From a theoretical perspective, gravitational waves emitted in extreme-mass-ratio inspirals encode detailed information about the spacetime geometry of supermassive black holes.


Such measurements allow for tests of whether astrophysical black holes are described by the Kerr solution of the Einstein equations.
Previous experience suggests that these observations are likely to confirm many predictions of general relativity, rule out certain modified gravity models, and constrain models of black hole formation.
However, the possibility of unexpected results remains.


Historically, opening new observational windows has led to unforeseen discoveries.
Gravitational-wave observations may provide insights into unresolved questions in gravitational physics, including the late-time acceleration of the Universe, galactic rotation curves, and the unification of quantum mechanics with gravity.


The gravitational wave revolution is entering its second phase.
The current detector network will continue to improve, eventually reaching design sensitivity and possibly being upgraded further.
But the most dramatic developments will come from new facilities:
\begin{itemize}
    \item \textbf{LISA (Laser Interferometer Space Antenna)}, formally adopted by ESA in January 2024, will launch in the mid-2030s to observe millihertz-frequency gravitational waves from space.
    LISA will detect supermassive black hole mergers, extreme mass ratio inspirals, and potentially stochastic backgrounds from the early universe—sources completely inaccessible to ground-based detectors due to seismic noise.


    \item \textbf{Einstein Telescope}, a proposed European underground facility with 10-kilometer arms, could improve sensitivity by a factor of 10 over current detectors~\cite{etscienceteam_einstein_2011}.
    Site selection is underway, with the Euregio Meuse-Rhine region (spanning Belgium, Netherlands, and Germany) and Sardinia, Italy, as leading candidates~\cite{gupta_characterizing_2024}.

    \item \textbf{Cosmic Explorer}, a U.S.\ proposal for two facilities with 40-kilometer and 20-kilometer arms, would achieve similar sensitivity improvements~\cite{hall_cosmic_2022}.
    The longer arms enable detection of binary black holes across cosmic time, potentially observing the first generation of black holes formed in the universe.

\end{itemize}


These next-generation facilities will not merely detect more events—they will enable qualitatively new science.
With sensitivities 10 times better than current detectors, they will observe every stellar-mass binary black hole merger in the observable universe, detect neutron star mergers out to cosmological distances, precisely test general relativity in regimes currently inaccessible, and potentially discover unexpected sources.
The computational challenges of handling tens of thousands of detections per year, extracting maximum information from each observation, and testing theoretical predictions with unprecedented precision will define much of the research landscape for the coming decades.


Additionally, pulsar timing arrays have recently reported evidence for a stochastic gravitational wave background at nanohertz frequencies~\cite{agazie_nanograv_2023,eptacollaborationandinptacollaboration:_second_2023,thefermi-latcollaboration*+_gammaray_2022}, likely from supermassive black hole binaries across the universe.
This opens yet another frequency window for gravitational wave observations, complementing ground-based and space-based detectors.


The gravitational wave revolution has moved from first detection to regular observations to planning for transformative new capabilities—all in less than a decade.
The field is simultaneously harvesting the scientific returns from current detections and preparing for facilities that will observe the gravitational wave sky in exquisite detail.
For researchers entering the field now, this creates a unique opportunity: contributing to both near-term challenges with current data and laying groundwork for the next generation of discoveries.

